\section{Conclusion}

\begin{guideline}
La conclusion est un bilan personnel et critique de votre stage. Elle
ne doit pas être un simple résumé de ce qui précède, mais une prise de
recul sincère sur votre expérience. Elle se compose de trois axes :
ce que vous retenez de la mission, votre regard critique sur ce qui
aurait pu être mieux, et vos perspectives professionnelles. Environ
1 à 2 pages.
\end{guideline}


\subsection{Bilan de la mission}

\begin{guideline}
Faites un bilan global de votre stage : avez-vous atteint les objectifs
fixés ? Quels sont les points forts de votre expérience ? Qu'est-ce que
ce stage vous a apporté sur le plan technique et humain ?
\end{guideline}

\instruction{[Dressez un bilan global honnête de votre stage. Répondez
à la question : « Qu'est-ce que j'ai accompli et appris ? » Soyez
factuel et précis. Faites le lien avec les objectifs fixés au départ
dans votre convention de stage.]}


\subsection{Retour critique}

\begin{guideline}
Prenez du recul sur votre mission : qu'est-ce qui n'a pas fonctionné
comme prévu ? Quelles décisions referiez-vous différemment ? Quelles
ont été vos principales difficultés et comment les avez-vous surmontées ?
Cette partie démontre votre maturité professionnelle.
\end{guideline}

\instruction{[Analysez avec honnêteté les limites de votre travail et
les difficultés rencontrées. Ex. : contraintes techniques imprévues,
délais non respectés, choix techniques que vous remettriez en question
aujourd'hui. Expliquez ce que vous en avez appris. C'est une marque
de maturité, pas une faiblesse.]}


\subsection{Perspectives professionnelles}

\begin{guideline}
Réfléchissez à la suite : qu'est-ce que ce stage a confirmé ou remis
en question dans vos aspirations ? Quelles compétences souhaitez-vous
développer en priorité ? Quel type de poste ou d'environnement
envisagez-vous ?
\end{guideline}

\instruction{[Expliquez comment ce stage a influencé votre vision de
votre futur métier d'ingénieur. Quelles compétences souhaitez-vous
approfondir ? Vers quel type de poste, d'entreprise ou de secteur
vous orientez-vous ? Soyez personnel et assumez vos choix.]}